\clearpage\section*{Correction EV1-T2}\label{EV1-T2-CORR}
\textbf{QCM :} 1 D, 2 C, 3 B, 4 D, 5 A, 6 A, 7 C, 8 A, 9 B, 10 C, 11 D, 12 C. Chaque réponse correcte : 0,5 point ; total 6.
\paragraph{Vérifications des automatismes.} Priorités ; somme de fractions ; soustraction des exposants ; coefficient multiplicateur 1,2 ; division par 0,8 ; produit 1,25 fois 0,8 ; résolution affine ; accroissement vertical divisé par l’horizontal ; moyenne ; angle divisé par 360 ; complément à 1 ; substitution entre parenthèses.
\subsection*{Exercice 1 -- Deux procédés de croissance}
\begin{enumerate}
\item A arithmétique, B géométrique ; $a_2=50,b_2=120$. \textit{Barème : 1 point. Résultat et justification élémentaire attendus ; demi-point possible pour démarche correcte inachevée.}
\item A : $\dfrac{4(30+60)}{2}=180$ ; B : $30(1+2+4+8)=450$. \textit{Barème : 2 points. Méthode ou justification : 1 ; calcul et conclusion : 1.}
\item $b_1=60>a_1=40$. Pour $n\geq1$, accroissement de B $30\times2^n\geq60$, supérieur à 10 ; chaque pas accroît l’écart. \textit{Barème : 2 points. Méthode ou justification : 1 ; calcul et conclusion : 1.}
\end{enumerate}
\subsection*{Exercice 2 -- Des racines imposées}
\begin{enumerate}
\item $f=a(x+1)(x-5)$ ; $-5a=10$, donc $a=-2$. \textit{Barème : 1 point. Résultat et justification élémentaire attendus ; demi-point possible pour démarche correcte inachevée.}
\item Coefficient négatif ; f positive entre les racines, négative ailleurs. Solutions $]-\infty;-1]\cup[5;+\infty[$. \textit{Barème : 2 points. Méthode ou justification : 1 ; calcul et conclusion : 1.}
\item $f=-2x^2+8x+10=-2(x-2)^2+18$. Maximum 18 en 2. \textit{Barème : 2 points. Méthode ou justification : 1 ; calcul et conclusion : 1.}
\end{enumerate}
\subsection*{Exercice 3 -- Une approximation contrôlée}
\begin{enumerate}
\item Taux $6+h$. \textit{Barème : 1 point. Résultat et justification élémentaire attendus ; demi-point possible pour démarche correcte inachevée.}
\item $f\prime(3)=6$, $f(3)=10$ ; $y=6x-8$. \textit{Barème : 1 point. Résultat et justification élémentaire attendus ; demi-point possible pour démarche correcte inachevée.}
\item Approximation 10,12 ; exact $10{,}1204$. Exact moins approché $0{,}0004=h^2$. \textit{Barème : 2 points. Méthode ou justification : 1 ; calcul et conclusion : 1.}
\end{enumerate}
\paragraph{Étalonnage.} Durée cible : 20–25 min pour les automatismes, 80–85 min pour les exercices, 10–15 min de relecture. Ces durées sont des estimations éditoriales à mesurer en classe. Accepter toute méthode correcte ; ne pas sanctionner deux fois une erreur numérique propagée avec cohérence. Le barème proposé ne vaut pas décision d’harmonisation de l’équipe.
